\section*{Bab \ref{ch:g10:reaction-progress-table} --- Tabel Kemajuan Reaksi}

\begin{solution}{exo:g10:reaction-progress-table:1}
\ce{H2}: $4.0 - 2x$; \ce{O2}: $3.0 - x$; \ce{H2O}: $2x$. Gas hidrogen
habis pada $x = 2.0$, gas oksigen pada $x = 3.0$: $x_{\max} = \qty{2.0}{mol}$,
hidrogen adalah pembatas. \omterm{def:g10:reaction-progress-table:system}{Keadaan akhir}: \qty{0}{mol} \ce{H2},
\qty{1.0}{mol} \ce{O2}, \qty{4.0}{mol} \ce{H2O}.
\end{solution}

\begin{solution}{exo:g10:reaction-progress-table:2}
Karbon akan habis pada $x = 3.0$, gas oksigen pada $x = 2.0$:
$x_{\max} = \qty{2.0}{mol}$, oksigen adalah pembatas. \omterm{def:g10:reaction-progress-table:system}{Keadaan akhir}:
\qty{1.0}{mol} karbon, tanpa oksigen, \qty{2.0}{mol} karbon dioksida.
\end{solution}

\begin{solution}{exo:g10:reaction-progress-table:3}
$x_{\max} = \qty{0.30}{mol}$ (belerang pembatas). \omterm{def:g10:reaction-progress-table:system}{Keadaan akhir}:
\qty{0.20}{mol} besi, tanpa belerang, \qty{0.30}{mol} besi sulfida.
\end{solution}

\begin{solution}{exo:g10:reaction-progress-table:4}
Gas nitrogen habis pada $x = 1.0$, gas hidrogen pada $x = 2.4/3 = 0.80$:
$x_{\max} = \qty{0.80}{mol}$. \omterm{def:g10:reaction-progress-table:system}{Keadaan akhir}: \qty{0.20}{mol} \ce{N2},
tanpa \ce{H2}, $2 \times 0.80 = \qty{1.6}{mol}$ \ce{NH3}.
\end{solution}

\begin{solution}{exo:g10:reaction-progress-table:5}
(b): $4/2 = 2/1$. Pada (a) karbon monoksida menjadi pembatas, pada (c)
juga.
\end{solution}

\begin{solution}{exo:g10:reaction-progress-table:6}
$n(\ce{CH4}) = 32 / 16.0 = \qty{2.0}{mol}$, jadi $x_{\max} = \qty{2.0}{mol}$
untuk \omterm{def:g10:reaction-progress-table:stoichiometric}{campuran stoikiometris}. Gas oksigen: $2 \times 2.0 = \qty{4.0}{mol}$,
yaitu $4.0 \times 32.0 = \qty{128}{g}$. Air:
$\qty{4.0}{mol} \times \qty{18.0}{g/mol} = \qty{72}{g}$.
\end{solution}

\begin{solution}{exo:g10:reaction-progress-table:7}
$n(\ce{C3H8}) = 1000 / 44.0 = \qty{22.7}{mol}$; karbon dioksida
$3 \times 22.7 = \qty{68.2}{mol}$; $V = 68.2 \times 24.1 =
\qty{1.64e3}{L}$, sekitar \qty{1.6}{m^3}.
\end{solution}

\begin{solution}{exo:g10:reaction-progress-table:8}
Pada $x = \qty{1.0}{mol}$: \qty{1.0}{mol} \ce{CH4}, \qty{1.0}{mol}
\ce{O2}, \qty{1.0}{mol} \ce{CO2}, \qty{2.0}{mol} \ce{H2O}. Gas oksigen
habis pada $x = \qty{1.5}{mol}$: reaksi berhenti di situ, sehingga
kemajuan yang lebih besar tidak pernah tercapai.
\end{solution}

\begin{solution}{exo:g10:reaction-progress-table:9}
$n_0(\ce{Zn}) = 1.30 / 65.4 = \qty{0.0199}{mol}$;
$n_0(\ce{Cu^{2+}}) = 0.10 \times 0.1000 = \qty{0.0100}{mol}$. \omterm{def:g9:ions:ion}{Ion-ion}
tembaga adalah pembatas: $x_{\max} = \qty{0.0100}{mol}$. Tembaga yang
terbentuk: $0.0100 \times 63.5 = \qty{0.635}{g}$; seng yang tersisa:
$(0.0199 - 0.0100) \times 65.4 = \qty{0.65}{g}$. Tidak ada \omterm{def:g9:ions:ion}{ion} tembaga
yang tersisa: warna birunya telah hilang.
\end{solution}

\begin{solution}{exo:g10:reaction-progress-table:10}
$n_0(\ce{Mg}) = 0.48 / 24.3 = \qty{0.0198}{mol}$, yang akan habis pada
$x = 0.0198$; $n_0(\ce{H+}) = \qty{0.020}{mol}$, yang akan habis pada
$x = 0.010$. Asamnya yang menjadi pembatas: $x_{\max} = \qty{0.010}{mol}$,
jadi \qty{0.010}{mol} gas hidrogen, $0.010 \times 24.1 = \qty{0.24}{L}$.
\end{solution}

\begin{solution}{exo:g10:reaction-progress-table:11}
\omterm{def:g10:reaction-progress-table:stoichiometric}{Campuran stoikiometris} dengan \qty{0.30}{mol} belerang memerlukan
\qty{0.30}{mol} besi; ada \qty{0.20}{mol} lebih banyak, yaitu
$0.20 / 0.30 \approx \qty{67}{\percent}$ berlebih.
\end{solution}

\begin{solution}{exo:g10:reaction-progress-table:12}
\omterm{def:g10:the-mole:amount}{Jumlah zatnya} sama: $m(\ce{Fe})/55.8 = m(\ce{S})/32.1$, dengan
$m(\ce{Fe}) + m(\ce{S}) = \qty{10.0}{g}$. Jadi
$m(\ce{Fe}) = 10.0 \times 55.8 / (55.8 + 32.1) = \qty{6.35}{g}$ dan
$m(\ce{S}) = \qty{3.65}{g}$.
\end{solution}

\begin{solution}{exo:g10:reaction-progress-table:13}
\begin{enumerate}
  \item $n_0(\ce{CaCO3}) = 10.0 / 100.1 = \qty{0.0999}{mol}$, satu-satunya
        \omterm{def:g7:chemical-reactions:reactant}{pereaksi}: $x_{\max} = \qty{0.0999}{mol}$; reaksi itu menghasilkan
        \qty{0.0999}{mol} \ce{CaO} dan \ce{CO2}, yaitu
        $0.0999 \times 24.1 = \qty{2.41}{L}$ karbon dioksida.
  \item $n_0(\ce{CaO}) = \qty{0.0999}{mol}$, $n_0(\ce{H2O}) = 1.8 / 18.0
        = \qty{0.100}{mol}$: kapur tohor menjadi pembatas (tipis sekali);
        \omterm{def:g6:pure-substances-and-mixtures:mixture}{campuran} itu hampir stoikiometris. \ce{Ca(OH)2}:
        $M = 40.1 + 2 \times (16.0 + 1.0) = \qty{74.1}{g/mol}$, massanya
        $0.0999 \times 74.1 = \qty{7.40}{g}$.
\end{enumerate}
\end{solution}

\begin{solution}{exo:g10:reaction-progress-table:14}
$n_0(\ce{H+}) = \qty{0.050}{mol}$ di setiap labu; \omterm{def:g9:ions:ion}{ion} ini akan habis
pada $x = \qty{0.025}{mol}$. Magnesium: 0.15, 0.30, 0.45, 0.60, 0.75,
\qty{0.90}{g} memberikan 0.0062, 0.0123, 0.0185, 0.0247, 0.0309,
\qty{0.0370}{mol}. Pada empat labu pertama magnesium adalah pembatas
dan gas hidrogennya $n(\ce{Mg}) \times 24.1$: 0.15, 0.30, 0.45, dan
\qty{0.60}{L}. Pada dua labu terakhir asamnya yang menjadi pembatas:
$0.025 \times 24.1
= \qty{0.60}{L}$. Grafiknya berupa garis lurus melalui titik asal
sampai sekitar \qty{0.61}{g} magnesium, lalu menjadi mendatar pada
\qty{0.60}{L}.
\end{solution}

\begin{solution}{exo:g10:reaction-progress-table:15}
$n_0(\text{etanol}) = 4.6 / 46.0 = \qty{0.100}{mol}$ (akan habis pada
$x = 0.100$); $n_0(\ce{O2}) = 4.8 / 24.1 = \qty{0.199}{mol}$ (akan habis
pada $x = 0.199/3 = 0.0664$). Gas oksigen adalah pembatas,
$x_{\max} = \qty{0.0664}{mol}$. Etanol yang tersisa: $0.100 - 0.0664 =
\qty{0.0336}{mol}$, yaitu $0.0336 \times 46.0 = \qty{1.5}{g}$. Karbon
dioksida: $2 \times 0.0664 = \qty{0.133}{mol}$, yaitu
$0.133 \times 24.1 = \qty{3.2}{L}$.
\end{solution}

\begin{solution}{pb:g10:reaction-progress-table:1}
\textbf{1.} Kiri: 2 Na, 6 N. Kanan: 2 Na, $3 \times 2 = 6$ N. Setara.

\textbf{2.} Kiri: 10 Na, 2 K, 2 N, 6 O. Kanan: K: 2; Na: $5 \times 2 =
10$; O: $1 + 5 = 6$; N: 2. Setara.

\textbf{3.} Bab tentang \omterm{def:g9:periodic-table-first-look:periodic-table}{tabel periodik} (\cref{ch:g9:periodic-table-first-look}):
natrium bereaksi hebat dengan air, menghasilkan larutan yang korosif dan
gas hidrogen. Setelah tabrakan, kantong itu bersentuhan dengan kulit,
mata, dan \omterm{def:g4:air-a-mixture-of-gases:air}{udara} lembap embusan napas.

\textbf{4.} $M(\ce{NaN3}) = 23.0 + 3 \times 14.0 = \qty{65.0}{g/mol}$;
$M(\ce{KNO3}) = 39.1 + 14.0 + 3 \times 16.0 = \qty{101.1}{g/mol}$.

\textbf{5.} \ce{NaN3}: $1.00 - 2x$; \ce{Na}: $2x$; \ce{N2}: $3x$.

\textbf{6.} $1.00 - 2x = 0$ pada $x_{\max} = \qty{0.500}{mol}$; natrium
azida, satu-satunya \omterm{def:g7:chemical-reactions:reactant}{pereaksi}, adalah pembatas.

\textbf{7.} \ce{Na}: $2 \times 0.500 = \qty{1.00}{mol}$; \ce{N2}:
$3 \times 0.500 = \qty{1.50}{mol}$.

\textbf{8.} $1.50 \times 24.1 = \qty{36.2}{L}$.

\textbf{9.} \ce{Na}: $1.00 - 10x$; \ce{KNO3}: $n - 2x$; \ce{K2O}: $x$;
\ce{Na2O}: $5x$; \ce{N2}: $x$.

\textbf{10.} $1.00/10 = n/2$, jadi $n = \qty{0.200}{mol}$, yaitu
$0.200 \times 101.1 = \qty{20.2}{g}$.

\textbf{11.} $x_{\max} = \qty{0.100}{mol}$: tanpa natrium, tanpa kalium
nitrat, \qty{0.100}{mol} \ce{K2O}, \qty{0.500}{mol} \ce{Na2O},
\qty{0.100}{mol} \ce{N2}.

\textbf{12.} Natrium akan habis pada $x = 0.100$, kalium nitrat pada
$x = 0.150/2 = 0.075$: nitratnya menjadi pembatas, $x_{\max} =
\qty{0.075}{mol}$, dan $1.00 - 10 \times 0.075 = \qty{0.25}{mol}$ \omterm{def:g9:periodic-table-first-look:metal}{logam}
natrium akan tersisa di dalam kantong, siap bereaksi dengan kelembapan.

\textbf{13.} $1.50 + 0.100 = \qty{1.60}{mol}$ nitrogen per \omterm{def:g10:the-mole:amount}{mol} natrium
azida.

\textbf{14.} $n = 60 / 24.1 = \qty{2.49}{mol}$.

\textbf{15.} $2.49 / 1.60 = \qty{1.56}{mol}$ natrium azida.

\textbf{16.} $1.56 \times 65.0 = \qty{101}{g}$.

\textbf{17.} $0.200 \times 1.56 = \qty{0.311}{mol}$, yaitu
$0.311 \times 101.1 = \qty{31.5}{g}$ kalium nitrat.

\textbf{18.} $2.49 / 1.50 = \qty{1.66}{mol}$, yaitu
$1.66 \times 65.0 = \qty{108}{g}$: reaksi kedua menghemat sekitar
\qty{7}{g} natrium azida.

\textbf{19.} \omterm{def:g10:reaction-progress-table:limiting}{Reaksi tuntas} hanya berhenti ketika \omterm{def:g10:reaction-progress-table:limiting}{pereaksi pembatasnya}
habis terpakai; di sini natrium azida adalah satu-satunya \omterm{def:g7:chemical-reactions:reactant}{pereaksi},
sehingga pada \omterm{def:g10:reaction-progress-table:system}{keadaan akhir} \omterm{def:g10:the-mole:amount}{jumlah zatnya} $1.00 - 2x_{\max} = 0$. Tidak
ada zat beracun yang tersisa di dalam kantong.

\textbf{20.} Sekitar \qty{101}{g} natrium azida mengisi kantong \omterm{def:g4:air-a-mixture-of-gases:air}{udara}
\qty{60}{L}.
\end{solution}
